\subsection{Резисторы}
%%%%%%%%%%%%%%%%%%%%%%%%%%%%%%%%%%%%%%%%%%%%%%%%%%%%
%%%%%%%%%%%%%%%%%%%%%%%%%%%%%%%%%%%%%%%%%%%%%%%%%%%%
\z{Если предположить, что число электронов проводимости в металле равно числу атомов, то какой будет средняя скорость электронов проводимости в серебряной проволоке диаметром $d = 1\,\text{мм}$, по которой идет ток $I = 30\,\text{А}$?}
%
{$v = \frac{4I\mu}{e\pi d^2 \rho N_A} \approx 4\,\text{мм/с}$}

%%%%%%%%%%%%%%%%%%%%%%%%%%%%%%%%%%%%%%%%%%%%%%%%%%%%
\z{Рассчитать сопротивление медного трамвайного провода длиной $l = 5\,\text{км}$, если площадь его сечения $S = 0.65\,\text{см}^2$.}
%
{$R = \rho \frac{l}{S} = 1.3\,\text{Ом}$}

%%%%%%%%%%%%%%%%%%%%%%%%%%%%%%%%%%%%%%%%%%%%%%%%%%%%
\z{После протягивания проволоки через волочильный станок длина ее увеличилась в $4$ раза. Каким стало сопротивление этой проволоки, если до обработки сопротивление было $R = 20\,\text{Ом}$.}
%
{$R_{\text{н}} = 16R = 320\,\text{Ом}$}

%%%%%%%%%%%%%%%%%%%%%%%%%%%%%%%%%%%%%%%%%%%%%%%%%%%%
\z[{\begin{wrapfigure}[5]{r}{0.3\textwidth}%
		\vspace{-20pt}%
		\begin{flushright}%
 			\includegraphics[width=0.3\textwidth]{3.2.4}%
		\end{flushright}%
	\end{wrapfigure}%
	}]%
{Четыре одинаковых резистора ученик включил в цепь, так как показано на рисунке. Какое он получит сопротивление между точками $A$ и $B$?}
%
{$\frac{R}{2}$}

%%%%%%%%%%%%%%%%%%%%%%%%%%%%%%%%%%%%%%%%%%%%%%%%%%%%
\z{Из проволоки сопротивлением $R = 25\,\text{Ом}$ сделано кольцо. Где следует присоединить провода, подводящие ток, чтобы сопротивление кольца было $r = 4\,\text{Ом}$? Длина проволоки $l$.}
%
{$\frac{x}{l} = \frac{1}{2}\left(1 - \sqrt{1 - \frac{4r}{R}}\right) = 0.2$}

%%%%%%%%%%%%%%%%%%%%%%%%%%%%%%%%%%%%%%%%%%%%%%%%%%%%
\z[{\begin{wrapfigure}[3]{r}{0.2\textwidth}%
		\vspace{-20pt}%
		\begin{flushright}%
 			\includegraphics[width=0.2\textwidth]{3.2.6}%
		\end{flushright}%
	\end{wrapfigure}%
	}]%
{Определить сопротивление $R_0$ проволочного каркаса в виде квадрата с диагоналями, спаянными в центре. Каркас включен в цепь точками $A$ и $B$. Сопротивление стороны квадрата $r$.}
%
{$(2 - \sqrt{2})R = 0.59R$}

%%%%%%%%%%%%%%%%%%%%%%%%%%%%%%%%%%%%%%%%%%%%%%%%%%%%
\z{Четыре одинаковых сопротивления, каждое из которых равно $r$, соединяют различными способами. Определить эквивалентное сопротивление.
		\begin{center}%
 			\includegraphics[width=0.8\textwidth]{3.2.7}%
		\end{center}%
}
%
{\puan $4R$;~\puan $\frac{R}{4}$;~\puan $R$;~\puan $R$;~\puan $\frac{4}{3}R$;~\puan $\frac{5}{2}R$;~\puan $\frac{5}{3}R$;~\puan $\frac{3}{5}R$;~\puan $\frac{3}{5}R$}
